% GENERATED FILE. Edit canonical lessons, then run npm run generate:books.
% canonical-source-hash: fnv1a64:88de8701598ea42f
% canonical-lessons: KA-C319-yuddha, KA-C319-sainika, KA-C319-akramisu, KA-C319-kalla, KA-C319-odihogu

\chapter{A Thief and a War}
\label{ch:ka-a-thief-and-a-war}
\hlchaptermodality{319}

\begin{chapteropening}
\textbf{By the end of this chapter:} \emph{I can say a war, a soldier, to invade, a thief and to run away.}

Complete the last lesson of chapter 319: I can say a war, a soldier, to invade, a thief and to run away.
\end{chapteropening}

\section[\texorpdfstring{\kn{ಯುದ್ಧ} (yuddha)}{ಯುದ್ಧ (yuddha)}]{\kn{ಯುದ್ಧ} (yuddha) — a war; a battle}

\label{lesson:KA-C319-yuddha}



\begin{quote}
Before the new one: say the Kannada for chalk, then the Kannada for a ruler.
\end{quote}

\subsection*{You'll want to know: \kn{ಯುದ್ಧ}}
\textbf{\kn{ಯುದ್ಧ}} — \emph{yuddha} — \textquotedblleft{}a war; a battle\textquotedblright{}.

\begin{quote}
\kn{ಯುದ್ಧದ} \kn{ಫಲಿತಾಂಶ} \kn{ಹವಾಮಾನದ} \kn{ಮೇಲೆ} \kn{ಅವಲಂಬಿಸಿತ್ತು}.
\end{quote}

\begin{quote}
\emph{yuddhada phalitāṃśa havāmānada mēle avalambisittu.}
\end{quote}

\begin{quote}
The outcome of the war depended on the weather.
\end{quote}

With \kn{ಮಾಡು} it makes a verb: \kn{ಯುದ್ಧ} \kn{ಮಾಡು}, to fight a war.

\begin{cousinweb}
\textbf{\kn{ಯುದ್ಧ}} (\emph{yuddha}) is borrowed from Sanskrit \emph{yuddha}, \textquotedblleft{}fought\textquotedblright{}, and as a noun \textquotedblleft{}a battle\textquotedblright{}, the past participle of the root \emph{yudh}, \textquotedblleft{}to fight\textquotedblright{} (CDIAL 10499). Kannada keeps the Sanskrit form whole, as it does many learned nouns.
\end{cousinweb}

\begin{grammarlens}[title={What you've built}]
One more word for telling what happened.
\end{grammarlens}

\subsection*{Your turn}
\begin{itemize}
\raggedright
  \item \emph{Say it:} \emph{yuddha}
  \item \emph{Say it:} \emph{yuddha}, once more
  \item \emph{Say it:} \emph{yuddhada phalitāṃśa havāmānada mēle avalambisittu.} — The outcome of the war depended on the weather.
  \item \emph{Recall:} say the Kannada for chalk, then the Kannada for a ruler, then say \emph{yuddha} again
\end{itemize}

\begin{tcolorbox}[breakable,colback=teal!4,colframe=teal!35!black,title={Before you move on}]
What does \kn{ಯುದ್ಧ} mean? (\textquotedblleft{}A war; a battle\textquotedblright{}.) Say it once more.
\end{tcolorbox}
\section[\texorpdfstring{\kn{ಸೈನಿಕ} (sainika)}{ಸೈನಿಕ (sainika)}]{\kn{ಸೈನಿಕ} (sainika) — a soldier}

\label{lesson:KA-C319-sainika}



\begin{quote}
Before the new one: say the Kannada for a ruler, then the Kannada for a war.
\end{quote}

\subsection*{You'll want to know: \kn{ಸೈನಿಕ}}
\textbf{\kn{ಸೈನಿಕ}} — \emph{sainika} — \textquotedblleft{}a soldier\textquotedblright{}.

\begin{quote}
\kn{ಸೈನಿಕರು} \kn{ಯುದ್ಧದಿಂದ} \kn{ಮನೆಗೆ} \kn{ಹಿಂತಿರುಗಿದರು}.
\end{quote}

\begin{quote}
\emph{sainikaru yuddhadinda manege hintirugidaru.}
\end{quote}

\begin{quote}
The soldiers came back home from the war.
\end{quote}

Plural: \kn{ಸೈನಿಕರು}, \emph{sainikaru}.

\begin{cousinweb}
\textbf{\kn{ಸೈನಿಕ}} (\emph{sainika}) is borrowed from Sanskrit \emph{sainika}, \textquotedblleft{}belonging to an army\textquotedblright{}, and as a noun \textquotedblleft{}a soldier\textquotedblright{} (Monier-Williams). It is built on \emph{senā}, \textquotedblleft{}an army\textquotedblright{}, with the ending -\emph{ika}, \textquotedblleft{}one who belongs to\textquotedblright{}.
\end{cousinweb}

\begin{grammarlens}[title={What you've built}]
One more word for telling what happened.
\end{grammarlens}

\subsection*{Your turn}
\begin{itemize}
\raggedright
  \item \emph{Say it:} \emph{sainika}
  \item \emph{Say it:} \emph{sainika}, once more
  \item \emph{Say it:} \emph{sainikaru yuddhadinda manege hintirugidaru.} — The soldiers came back home from the war.
  \item \emph{Recall:} say the Kannada for a ruler, then the Kannada for a war, then say \emph{sainika} again
\end{itemize}

\begin{tcolorbox}[breakable,colback=teal!4,colframe=teal!35!black,title={Before you move on}]
What does \kn{ಸೈನಿಕ} mean? (\textquotedblleft{}A soldier\textquotedblright{}.) Say it once more.
\end{tcolorbox}
\section[\texorpdfstring{\kn{ಆಕ್ರಮಿಸು} (ākramisu)}{ಆಕ್ರಮಿಸು (ākramisu)}]{\kn{ಆಕ್ರಮಿಸು} (ākramisu) — to invade; to attack}

\label{lesson:KA-C319-akramisu}



\begin{quote}
Before the new one: say the Kannada for a war, then the Kannada for a soldier.
\end{quote}

\subsection*{You'll want to know: \kn{ಆಕ್ರಮಿಸು}}
\textbf{\kn{ಆಕ್ರಮಿಸು}} — \emph{ākramisu} — \textquotedblleft{}to invade; to attack\textquotedblright{}.

\begin{quote}
\kn{ಸೈನಿಕರು} \kn{ಆ} \kn{ದೇಶವನ್ನು} \kn{ಆಕ್ರಮಿಸಿದರು}.
\end{quote}

\begin{quote}
\emph{sainikaru ā dēśavannu ākramisidaru.}
\end{quote}

\begin{quote}
The soldiers invaded that country.
\end{quote}

Past (I …): \kn{ಆಕ್ರಮಿಸಿದೆ}, \emph{ākramiside}.

What is attacked takes -\kn{ಅನ್ನು}: \kn{ದೇಶವನ್ನು} \kn{ಆಕ್ರಮಿಸು}, to invade a country.

\begin{cousinweb}
\textbf{\kn{ಆಕ್ರಮಿಸು}} (\emph{ākramisu}) is built from Sanskrit \emph{ākramaṇa}, \textquotedblleft{}stepping on, an advance\textquotedblright{} (CDIAL 1016), from \emph{ā-} \textquotedblleft{}towards\textquotedblright{} and \emph{kram} \textquotedblleft{}to step\textquotedblright{}, with the Kannada ending -\emph{isu} that turns a Sanskrit noun into a verb.
\end{cousinweb}

\begin{grammarlens}[title={What you've built}]
One more word for telling what happened.
\end{grammarlens}

\subsection*{Your turn}
\begin{itemize}
\raggedright
  \item \emph{Say it:} \emph{ākramisu}
  \item \emph{Say it:} \emph{ākramisu}, once more
  \item \emph{Say it:} \emph{sainikaru ā dēśavannu ākramisidaru.} — The soldiers invaded that country.
  \item \emph{Recall:} say the Kannada for a war, then the Kannada for a soldier, then say \emph{ākramisu} again
\end{itemize}

\begin{tcolorbox}[breakable,colback=teal!4,colframe=teal!35!black,title={Before you move on}]
What does \kn{ಆಕ್ರಮಿಸು} mean? (\textquotedblleft{}To invade; to attack\textquotedblright{}.) Say it once more.
\end{tcolorbox}
\section[\texorpdfstring{\kn{ಕಳ್ಳ} (kaḷḷa)}{ಕಳ್ಳ (kaḷḷa)}]{\kn{ಕಳ್ಳ} (kaḷḷa) — a thief}

\label{lesson:KA-C319-kalla}



\begin{quote}
Before the new one: say the Kannada for a soldier, then the Kannada for to invade.
\end{quote}

\subsection*{You'll want to know: \kn{ಕಳ್ಳ}}
\textbf{\kn{ಕಳ್ಳ}} — \emph{kaḷḷa} — \textquotedblleft{}a thief\textquotedblright{}.

\begin{quote}
\kn{ಕಳ್ಳ} \kn{ನನ್ನ} \kn{ಚೀಲವನ್ನು} \kn{ಕದ್ದನು}.
\end{quote}

\begin{quote}
\emph{kaḷḷa nanna cīlavannu kaddanu.}
\end{quote}

\begin{quote}
The thief stole my bag.
\end{quote}

Plural: \kn{ಕಳ್ಳರು}, \emph{kaḷḷaru}.

A woman thief is \kn{ಕಳ್ಳಿ}. The verb beside it is \kn{ಕದಿ}, to steal: \kn{ಕದ್ದನು}, he stole.

\begin{cousinweb}
\textbf{\kn{ಕಳ್ಳ}} (\emph{kaḷḷa}) is native Dravidian: DEDR 1372 lists Kannada \emph{kaḷḷa} \textquotedblleft{}a thief\textquotedblright{} beside Tamil \emph{kaḷḷaṉ} \textquotedblleft{}a thief, a robber\textquotedblright{} and \emph{kaḷ} \textquotedblleft{}to steal\textquotedblright{}.
\end{cousinweb}

\begin{grammarlens}[title={What you've built}]
One more word for telling what happened.
\end{grammarlens}

\subsection*{Your turn}
\begin{itemize}
\raggedright
  \item \emph{Say it:} \emph{kaḷḷa}
  \item \emph{Say it:} \emph{kaḷḷa}, once more
  \item \emph{Say it:} \emph{kaḷḷa nanna cīlavannu kaddanu.} — The thief stole my bag.
  \item \emph{Recall:} say the Kannada for a soldier, then the Kannada for to invade, then say \emph{kaḷḷa} again
\end{itemize}

\begin{tcolorbox}[breakable,colback=teal!4,colframe=teal!35!black,title={Before you move on}]
What does \kn{ಕಳ್ಳ} mean? (\textquotedblleft{}A thief\textquotedblright{}.) Say it once more.
\end{tcolorbox}
\section[\texorpdfstring{\kn{ಓಡಿಹೋಗು} (ōḍihōgu)}{ಓಡಿಹೋಗು (ōḍihōgu)}]{\kn{ಓಡಿಹೋಗು} (ōḍihōgu) — to run away; to escape}

\label{lesson:KA-C319-odihogu}



\begin{quote}
Before the new one: say the Kannada for to invade, then the Kannada for a thief.
\end{quote}

\subsection*{You'll want to know: \kn{ಓಡಿಹೋಗು}}
\textbf{\kn{ಓಡಿಹೋಗು}} — \emph{ōḍihōgu} — \textquotedblleft{}to run away; to escape\textquotedblright{}.

\begin{quote}
\kn{ಪೊಲೀಸರನ್ನು} \kn{ನೋಡಿ} \kn{ಕಳ್ಳ} \kn{ಓಡಿಹೋದನು}.
\end{quote}

\begin{quote}
\emph{polīsarannu nōḍi kaḷḷa ōḍihōdanu.}
\end{quote}

\begin{quote}
Seeing the police, the thief ran away.
\end{quote}

Past (I …): \kn{ಓಡಿಹೋದೆ}, \emph{ōḍihōde}.

\kn{ಓಡಿ} is \kn{ಓಡು}, to run, as \textquotedblleft{}having run\textquotedblright{}; \kn{ಹೋಗು} adds \textquotedblleft{}away\textquotedblright{}. Run and go make one verb.

\begin{cousinweb}
\textbf{\kn{ಓಡಿಹೋಗು}} (\emph{ōḍihōgu}) is a native Dravidian compound: \emph{ōḍu} \textquotedblleft{}to run\textquotedblright{} (DEDR 1041) and \emph{hōgu} \textquotedblleft{}to go\textquotedblright{} (DEDR 4572).
\end{cousinweb}

\begin{grammarlens}[title={What you've built}]
One more word for telling what happened.
\end{grammarlens}

\subsection*{Your turn}
\begin{itemize}
\raggedright
  \item \emph{Say it:} \emph{ōḍihōgu}
  \item \emph{Say it:} \emph{ōḍihōgu}, once more
  \item \emph{Say it:} \emph{polīsarannu nōḍi kaḷḷa ōḍihōdanu.} — Seeing the police, the thief ran away.
  \item \emph{Recall:} say the Kannada for to invade, then the Kannada for a thief, then say \emph{ōḍihōgu} again
\end{itemize}

\begin{tcolorbox}[breakable,colback=teal!4,colframe=teal!35!black,title={Before you move on}]
What does \kn{ಓಡಿಹೋಗು} mean? (\textquotedblleft{}To run away; to escape\textquotedblright{}.) Say it once more.
\end{tcolorbox}
